% =====================================================================
\chapter{Marco teórico}
\label{cap:marco}

Que «las noticias mueven la bolsa» mezcla dos ideas: que el precio cambia cuando llega
información nueva, y que los medios son el vehículo de esa información. Pero los medios también
informan sobre el propio mercado: una nota que cuenta que «el Merval subió 5\%» tiene sentimiento
positivo, y su causa es el precio, no al revés. Este capítulo presenta los conceptos con los que
se analiza esa doble relación y las herramientas de ciencia de datos con las que se la pone a
prueba.

% ---------------------------------------------------------------------
\section{Noticias, sentimiento y mercado}

\textbf{Shock informativo y eficiencia.} Un \textit{shock informativo} es la llegada de
información que se aparta de lo esperado. La hipótesis de mercados eficientes \citep{fama1970}
sostiene que la información pública se incorpora al precio tan rápido que no sirve para
anticiparlo; la evidencia para Argentina y la región es mixta \citep{zablotsky2002}. Con las
noticias se puede medir la \textit{atención}, es decir, cuánto se habla \citep{da2011,fang2009}, o
el \textit{contenido}, es decir, cómo se habla \citep{tetlock2007,ardia2022}. Esta tesina mide el
contenido, por una razón práctica que se explica en el capítulo~\ref{cap:datos}: la cantidad de
notas recolectadas depende de qué capturó un archivo web, no solo de cuánto publicaron los medios.

\textbf{Sentimiento de los medios.} El sentimiento de una nota es el signo de lo que comunica:
positivo, negativo o neutro. Un sentimiento muy pesimista anticipa bajas transitorias
\citep{tetlock2007}, sobre todo en las recesiones \citep{garcia2013}; pero el sentimiento también
puede ser un \textit{reflejo} del mercado \citep{verma2026}. La tabla~\ref{tab:mecanismos} resume
cómo se reconoce cada dirección en los datos.

\begin{table}[H]
\centering\small
\begin{tabularx}{\linewidth}{@{}L{3.2cm}YY@{}}
\toprule
 & \textbf{Noticias \flecha{} mercado} & \textbf{Mercado \flecha{} noticias} \\
\midrule
Mecanismo & La nota trae información nueva o mueve el ánimo de los inversores & La nota cuenta lo que el mercado ya hizo \\
Qué se espera ver & El sentimiento de hoy anticipa el retorno de mañana & El retorno de hoy anticipa el sentimiento de mañana \\
En el modelo predictivo de esta tesina & Agregar el sentimiento mejora la predicción de si la acción sube & Agregar el retorno mejora la predicción del sentimiento \\
\bottomrule
\end{tabularx}
\caption{Las dos direcciones de la relación entre noticias y mercado y cómo se reconoce cada
una. Fuente: elaboración propia.}
\label{tab:mecanismos}
\end{table}

\textbf{Volatilidad.} La volatilidad mide cuánto se mueve un precio. En el S\&P~Merval se agrupa
en períodos de calma y de estrés \citep{segovia2022}, y las caídas pesan más que las subas
\citep{otuki2008}. El VIX de Chicago la mide con precios de opciones, como la volatilidad esperada
de los próximos 30 días \citep{cboe_vix}; en Argentina no hay un mercado de opciones con esa
liquidez, por lo que se usa la volatilidad realizada (sección~\ref{sec:vix}).

% ---------------------------------------------------------------------
\section{Aprendizaje supervisado para probar la relación}
\label{sec:ml}

La pregunta «¿las noticias ayudan a predecir al mercado?» se puede convertir en un problema de
\textbf{clasificación supervisada}: se le da a un modelo los datos de hoy y tiene que predecir una
clase de mañana (por ejemplo, «sube» o «no sube»). Si un modelo con noticias predice mejor que uno
sin noticias en datos que no vio, las noticias contienen información útil sobre el mercado.

\textbf{Regresión logística.} Estima la probabilidad de la clase 1 como una función de una
combinación lineal de las variables \citep{hosmer2013}:
\begin{equation}
P(y = 1 \mid x) = \frac{1}{1 + e^{-(\beta_0 + \beta_1 x_1 + \dots + \beta_k x_k)}}.
\end{equation}
Con las variables estandarizadas, el tamaño de cada coeficiente $\beta$ indica cuánto pesa esa
variable, y su signo, en qué dirección. Es el modelo principal porque es interpretable.

\textbf{Random forest.} Combina cientos de árboles de decisión, cada uno entrenado con una muestra
distinta de los datos y de las variables \citep{breiman2001}. Capta relaciones no lineales e
interacciones que la regresión logística no ve. Se usa como confirmación: si los dos modelos dicen
lo mismo, el resultado no depende de la forma del modelo.

\textbf{Evaluación con datos no vistos.} Un modelo se evalúa en datos que no usó para entrenar. Con
series de tiempo, la división tiene que respetar el orden temporal: entrenar con el pasado y
evaluar con el futuro. Si no, información del futuro se filtra al entrenamiento y las métricas
salen infladas, el problema de \textit{fuga de información} \citep{kapoor2023leakage}, que con
noticias cambia mucho el resultado \citep{xue2026temporal}.

\textbf{Métricas.} Con una predicción binaria se arma la \textbf{matriz de confusión}: cuántos
casos de cada clase real se predijeron en cada clase. De ella salen:
\begin{itemize}
  \item \textbf{Exactitud} (\textit{accuracy}): proporción de aciertos. Engaña si una clase es mucho
        más frecuente que la otra.
  \item \textbf{Balanced accuracy}: promedio de la proporción de aciertos dentro de cada clase real.
        Un modelo que predice siempre lo mismo obtiene 0,5.
  \item \textbf{F1}: media armónica entre la precisión (de lo que el modelo predijo como clase 1,
        cuánto lo era) y la sensibilidad (de lo que era clase 1, cuánto encontró). El
        \textbf{F1 macro} es el promedio simple del F1 de cada clase.
  \item \textbf{AUC}: área bajo la curva ROC \citep{fawcett2006}. Es la probabilidad de que el
        modelo le asigne más probabilidad a un caso de la clase 1 que a uno de la clase 0. Vale 0,5
        si el modelo no distingue y 1 si distingue perfecto. No depende del umbral de 0,5, por lo
        que es la métrica principal para comparar modelos.
\end{itemize}

\textbf{Intervalo de confianza con bootstrap.} Para saber si la diferencia de AUC entre dos modelos
es real o del azar, se remuestrea el test muchas veces y se mira cómo varía la diferencia. Con
series de tiempo se remuestrean \textit{bloques} de días consecutivos, para conservar la
dependencia entre días cercanos \citep{kunsch1989}.

% ---------------------------------------------------------------------
\section{Verificación clásica: VAR y causalidad de Granger}

Con datos observacionales no se prueba una causa, pero sí qué variable se mueve antes. $X$
\textit{causa en sentido de Granger} a $Y$ si el pasado de $X$ mejora la predicción de $Y$ más allá
de su propio pasado \citep{granger1969}. La prueba se hace dentro de un vector autorregresivo
(VAR), que explica cada variable con los rezagos de todas \citep{sims1980,lutkepohl2005}. Es la
misma idea del modelo predictivo, pero dentro de la muestra y con una prueba F: por eso se usa
como verificación del resultado principal.
